\section{Background}
\label{sec:background}

\subsection{Coding-agent evaluation and its artifacts}

A coding agent takes a repository and a natural-language maintenance task -- typically a
GitHub issue~\cite{SweAgentAgent2024,DemystifyingLlmBased2025} -- and produces a change to
that repository, judged by whether the project's
tests pass afterwards. Evaluations of such agents report aggregate rates over a fixed task
set: a percentage resolved, a proportion of empty patches, a best-of-$k$ curve.

Each reported rate rests on a chain of per-task records: which tasks were in scope, what the
agent produced for each, how each output was evaluated, and how outcomes that were neither
success nor failure -- timeouts, empty generations, runs that never started -- were treated.
The rate is a summary of that chain. Recovering the rate from released evidence requires the
chain, not the summary.

\subsection{What availability means for a run-level claim}

Artifact availability is usually treated as a repository link. For a run-level claim it is
finer. Each required element of the chain is a separate item with its own state, and a study
can be complete in one and empty in another: released patches with no evaluation outcomes,
released outcomes with no terminal classification, released everything with no statement of
which release the paper's table reports.

We use a six-state item vocabulary under a precedence, assigned per required item rather
than per study. Two states describe retrieved bytes: one where an independently computed
digest matches a recorded immutable digest under the same algorithm and canonical scope, one
where the complete bytes arrived but nothing verifies them. Four describe distinct reasons
for absence -- access is restricted, release is declared not to have happened, a locator was
attempted and returned no complete closure, or no locator for the item exists at all. The
last two are the ones most easily confused, and they are not the same: one is a study whose
artifact cannot be pulled, the other a study that published nothing to pull.

The distinctions are not bookkeeping. Collapsing \emph{attempted and incomplete} into
\emph{never published} would let a paywall read as a finding, and collapsing \emph{complete
bytes} into \emph{verified bytes} would let an unpinned download read as an identity check.

\subsection{Reproduction, transfer, and what blocks each}

Two questions are separable. Source reproduction asks whether a published comparison can be
re-derived under its own captured protocol. Protocol transfer asks whether a different,
compatible protocol applied to the same raw evidence would change the comparison's outcome.

Transfer needs the source's raw per-task evidence and a binding from that evidence to the
claim it supports, because the transferred estimate is projected back through the source's
own decision semantics. A source-side gap in a reported rule does not by itself end the
transfer question -- a frozen completion can become a new target protocol -- but a missing
binding does, because there is then no established mapping from released evidence to the
claim being transferred.

This paper does not evaluate whether transfers reverse published conclusions. The tier-two
sample was drawn, but its compatibility analysis was not executed, so transfer applicability
and reversal were not estimated.
